% ============================================
% DOC-27
% Extraído de main.pandoc.tex
% ============================================
\section*{\textbf{DOC-27} --- A Origem (1853)}
\markboth{DOC-27 --- A Origem (1853)}{}

\noindent \textit{Documento extraído do arquivo pericial.}

\vspace{0.5em}\noindent\rule{\textwidth}{0.5pt}\vspace{0.5em}

\noindent {\footnotesize\texttt{ARQUIVO P\'UBLICO DO CEAR\'A $\cdot$ DOCUMENTO AVULSO $\cdot$ S\'EC. XIX}}\\[-0.2em]

\noindent\textbf{Documento:} carta do Padre José Inácio de Sousa ao Presidente da Província do Ceará, Comissão de Socorros Públicos.\\
\noindent\textbf{Data:} 14 de março de 1853.\\
\noindent\textbf{Origem:} Freguesia de Nossa Senhora dos Remédios, Rio Quixeramobim.\\
\noindent\textbf{Estado de conservação:} papel manchado por umidade; a borda inferior apresenta uma substância escura, não identificada, que atravessou todas as folhas do maço. Análise química inconclusiva.\\
\noindent\textbf{Anexado por:} Dra. Lúcia Mendes, em 12 de fevereiro de 2026, com a anotação: \textquotedblleft{\textit{Isto prova que o diário não é o começo. O diário é só o eco.\textquotedblright{}}}

\noindent Excelentíssimo Senhor Presidente da Província,

\noindent Tomo a liberdade de escrever-lhe porque os vivos já não me ouvem e os mortos são muitos para enterrar sozinho.

\noindent Estou aqui desde novembro do ano passado, destacado para administrar os sacramentos aos retirantes que chegam à margem do rio. Recebi duzentas almas em janeiro. Duzentas e quarenta em fevereiro. Neste mês de março, já perdi a conta. Chegam aos magotes, famílias inteiras reduzidas a pó ambulante. Os olhos deles são a primeira coisa que morre. Antes do corpo. Os olhos perdem a luz, a água, a vida. Ficam abertos, mas não veem. Como se a seca tivesse entrado por eles e secado a alma por dentro.

\noindent Os que morrem são levados para uma vala a meia légua da freguesia. Não tenho cemitério consagrado. Não tenho madeira para caixões. Não tenho pano para mortalhas. Tenho apenas uma cova comprida, aberta na terra seca, onde os corpos são deitados em camadas. Terra sobre eles. Mais corpos. Mais terra. Os urubus aprenderam o caminho. Fazem um círculo no céu que não se desfaz. Nem de noite. Ficam empoleirados nas árvores ao redor, esperando.

\noindent E é sobre isto que preciso escrever-lhe.

\noindent Há três noites, não consigo rezar o ofício das almas sem sentir uma presença. Não uma presença espiritual, como as que estou acostumado a discernir em trinta anos de sacerdócio. É uma presença física. Ela chega com o cheiro.

\noindent O cheiro vem da vala.

\noindent Não é o cheiro comum da morte, que eu conheço bem. É um cheiro doce. Mais doce que incenso, mais doce que flor de mandacaru. Um cheiro que entra pelas narinas e se instala na garganta como um líquido grosso. Os retirantes sentem também. Eles se viram na direção da vala e ficam olhando. Não choram. Não rezam. Apenas olham. Como se esperassem que algo fosse sair de lá.

\noindent Eu mesmo fui até a vala na segunda noite.

\noindent Levei uma lanterna, água benta, um crucifixo. Coisas de padre. Coisas que sempre funcionaram. Cheguei à borda e iluminei o fundo.

\noindent O que vi não eram corpos.

\noindent Eram restos. Braços e pernas saindo da terra como raízes. Mãos abertas para cima, como se pedissem. Uma mão de criança, pequena, com os dedos retorcidos, apontava para o céu noturno como se acusasse alguém.

\noindent O cheiro era mais forte ali.

\noindent Ajoelhei para benzer a vala. Comecei o \textit{Requiem Aeternam}. Na terceira palavra, senti que a terra vibrou.

\noindent Não um tremor. Uma pulsação. Como se o chão tivesse um coração.

\noindent Ajoelhei. Pressionei a mão direita contra a terra. Ela estava morna. Não quente do sol — morna de dentro. A temperatura do corpo de um doente. A terra estava febril.

\noindent Enterrei a mão. Até o punho. A terra era macia, úmida — coisa impossível neste sertão seco. Senti algo roçar meus dedos. Algo liso. Mole. Algo que se moveu ao toque, como se estivesse vivo.

\noindent Tirei a mão.

\noindent A terra fechou sozinha sobre o buraco.

\noindent Um homem que estava perto — um retirante chamado Vicente, que perdeu a mulher e três filhos na semana passada — olhou para mim e disse:

\noindent \textquotedblleft{\textit{Seu padre, a vala tá com fome.\textquotedblright{}}}

\noindent Excelência, sei o que Vossa Senhoria pensará ao ler isto. Que o isolamento e a fome abalaram meu juízo. Que o calor e a mortandade turvaram minha razão. Não sei se é verdade. Talvez seja. Mas preciso registrar o que vi, para que, se eu morrer aqui — e acho que vou morrer aqui —, alguém saiba que esta vala não é uma vala comum.

\noindent No terceiro dia — ontem — acordei com a palma da mão direita marcada.

\noindent Uma marca escura, arredondada, no centro da palma. Parece uma queimadura, mas não dói. A pele está íntegra. A marca apenas... está lá. Como se alguém tivesse pressionado um selo quente contra minha carne enquanto eu dormia.

\noindent Tentei lavar. Não sai.

\noindent Tentei rezar. A marca não some.

\noindent Talvez seja só a terra que entrou na minha pele. Talvez seja a úlcera que está começando. Talvez seja o juízo que se vai.

\noindent Mas, esta noite, quando fechei os olhos, ouvi uma voz.

\noindent Não uma voz que vem de fora. Uma voz que vem de \textit{dentro da terra}. Ela repetia a mesma palavra, sem parar, como quem aprende a falar:

\noindent \textquotedblleft{\textit{Fome.\textquotedblright{}}}

\noindent \textquotedblleft{\textit{Fome.\textquotedblright{}}}

\noindent \textquotedblleft{\textit{Fome.\textquotedblright{}}}

\noindent A voz de uma criança. Mas nenhuma criança sobreviveu para falar.

\noindent Peço licença para encerrar esta carta. A lanterna está baixando.

\noindent Que Deus tenha misericórdia de nós.

\noindent \textit{Padre José Inácio de Sousa}\\
\noindent \textit{Freguesia de Nossa Senhora dos Remédios}\\
\noindent \textit{14 de março de 1853}
\noindent \textit{Nota do Arquivo Público do Ceará: a carta do Padre Sousa foi encontrada em 1916, dentro de um livro de tombo da freguesia, durante a reforma do prédio. O sacerdote faleceu em 15 de março de 1853 — um dia após redigir esta carta. A causa da morte foi registrada como \textquotedblleft{}fome voluntária\textquotedblright{}: o padre teria deixado de se alimentar, cedendo sua ração aos retirantes. O corpo foi enterrado na mesma vala que ele descreve.}
\noindent Anexa-se a este documento a ficha de catalogação de 1916, onde um funcionário anônimo escreveu a lápis, na margem: \textquotedblleft{\textit{O cheiro doce nunca saiu daqui. Transferir para o Arquivo Público Mineiro. Processo 1.260.\textquotedblright{} O funcionário não foi identificado. A anotação está na mesma caligrafia trêmula que aparece nas bordas do diário de Joaquim.}

\noindent \textbf{Julho de 2026 --- Belo Horizonte}\\
\noindent \textbf{Último registro profissional de Dra. Lúcia Mendes}

\noindent \textbf{Documento interno --- Autoencaminhamento para afastamento clínico}\\
\noindent \textbf{Profissional:} Lúcia Mendes, CRM-MG 28.391. \textbf{Data:} 28 de julho de 2026.

\noindent \textbf{Observação ética:}

\noindent ste registro não constitui laudo psicológico válido. Foi redigido para justificar afastamento imediato, interrupção de pesquisa e encaminhamento a avaliação independente. A primeira pessoa foi mantida por se tratar de anotação de urgência.

\noindent \textbf{Histórico resumido:}
\noindent Em janeiro de 2026, a profissional teve contato com o diário atribuído ao paciente 1.260, Joaquim Antônio Correia, no contexto de pesquisa em psicologia forense histórica. O material permaneceu sob guarda intermitente, com falhas posteriores de lacre e registro.

\noindent Nas semanas seguintes, foram observados sinais persistentes após manuseio do arquivo:

\begin{itemize}
\item alteração subjetiva de coloração no olho direito, ainda sem exame oftalmológico conclusivo;
\item percepção recorrente de odor adocicado sem fonte ambiental identificada;
\item sonhos repetitivos com vala, enfermaria e formulário em branco;
\item hipervigilância noturna e dificuldade de concentração;
\item impulso de retornar ao diário mesmo após decisão formal de encerramento da análise.
\end{itemize}

\noindent \textbf{Hipóteses de trabalho:}

\begin{itemize}
\item reação aguda ao material traumático pesquisado, com imagens intrusivas e insônia;
\item sugestão intensificada pela longa exposição documental;
\item identificação contratransferencial com o paciente 1.260;
\item fenômeno dissociativo leve, a confirmar por avaliação externa.
\end{itemize}

\noindent \textbf{Medidas adotadas:}

\begin{itemize}
\item suspensão de entrevistas, escrita acadêmica e circulação de cópias;
\item encaminhamento do original para custódia judicial;
\item solicitação de avaliação clínica por profissional sem contato prévio com o caso;
\item recomendação de que o diário não seja lido individualmente.
\end{itemize}

\noindent \textbf{Anotação não protocolada:}
\noindent A cópia impressa deste documento apresenta, ao final, a numeração ``1.262'' escrita a lápis ao lado do nome da profissional. O traço não corresponde à assinatura principal. Abaixo dele há uma linha vazia, sem identificação.

\noindent \textit{Lúcia Mendes --- CRM-MG 28.391}\\
\noindent Registro encerrado às 3h14, conforme horário anotado no rodapé da cópia.

\vspace{0.5em}
\noindent\textit{\textbf{↪ Rotas:} \rota{DOC-01} \rotasep \rota{LUC-02} \rotasep \rota{EPI-01}}

}